\section{The Hadamard Matrix and Its Operator Norms}

For clarity, write $H_m$ for the Sylvester--Hadamard matrix of order $m=2^k$. Thus $H_1=(1)$, and for $m=2n$,
\[
H_m = H_{2n} = \begin{pmatrix} H_n & H_n\\ H_n & -H_n \end{pmatrix}.
\]

We prove the following properties.

\subsection*{(a) $H$ is symmetric and has only $\pm 1$ entries}

We prove by induction on $k$ that $H_{2^k}$ is symmetric and every entry of $H_{2^k}$ is equal to $+1$ or $-1$.

For $k=0$, we have $H_1=(1)$, which is symmetric and has only $+1$ entries.

Assume now that $H_n$ is symmetric and has only $\pm 1$ entries. Consider
\[
H_{2n} = \begin{pmatrix} H_n & H_n\\ H_n & -H_n \end{pmatrix}.
\]
Every block $H_n$ and $-H_n$ has only $\pm 1$ entries, so $H_{2n}$ also has only $\pm 1$ entries. Moreover, since $H_n^\top=H_n$, we have
\[
H_{2n}^\top = \begin{pmatrix} H_n^\top & H_n^\top\\ H_n^\top & (-H_n)^\top \end{pmatrix} = \begin{pmatrix} H_n & H_n\\ H_n & -H_n \end{pmatrix} = H_{2n}.
\]
Thus $H_{2n}$ is symmetric. By induction, $H_m$ is symmetric and has only $\pm 1$ entries for every $m=2^k$.

\subsection*{(b) $H^\top H=H^2=mI_m$}

Since $H_m$ is symmetric by part~(a), it suffices to prove $H_m^2=mI_m$.

We prove this by induction on $m=2^k$. For $m=1$, $H_1^2=(1)^2=(1)=I_1$, so the claim holds.

Assume now that for some $n=2^{k-1}$, we have $H_n^2=nI_n$. Let $m=2n$. Then
\[
H_{2n}^2 = \begin{pmatrix} H_n & H_n\\ H_n & -H_n \end{pmatrix} \begin{pmatrix} H_n & H_n\\ H_n & -H_n \end{pmatrix} = \begin{pmatrix} 2H_n^2 & 0\\ 0 & 2H_n^2 \end{pmatrix}.
\]
Using the induction hypothesis $H_n^2=nI_n$, we get
\[
H_{2n}^2 = \begin{pmatrix} 2nI_n & 0\\ 0 & 2nI_n \end{pmatrix} = 2nI_{2n}.
\]
Since $m=2n$, this is exactly $H_m^2=mI_m$. Thus by induction $H^2=mI_m$, and since $H^\top=H$, we have $H^\top H=H^2=mI_m$.

\subsection*{(c) Operator norms of $H$}

We use the following well-known characterizations of induced operator norms. For any $A\in\mathbb{R}^{m\times m}$:
\begin{itemize}
\item \textbf{Spectral norm:} $\|A\|_{2\to 2} = \sqrt{\lambda_{\max}(A^\top A)}$, where $\lambda_{\max}$ denotes the largest eigenvalue. \textit{Proof:} For any $x\in\mathbb{R}^m$, $\|Ax\|_2^2 = x^\top A^\top A x \leq \lambda_{\max}(A^\top A)\|x\|_2^2$, with equality when $x$ is the corresponding eigenvector.
\item \textbf{$\ell^1$ norm:} $\|A\|_{1\to 1} = \max_{1\leq j\leq m}\sum_{i=1}^m |a_{ij}|$ (maximum absolute column sum). \textit{Proof:} For $x$ with $\|x\|_1=1$, $\|Ax\|_1 = \sum_i|\sum_j a_{ij}x_j| \leq \sum_i\sum_j|a_{ij}||x_j| = \sum_j|x_j|\sum_i|a_{ij}| \leq \max_j\sum_i|a_{ij}|$, with equality at $x=e_{j^*}$ for the maximizing column $j^*$.
\item \textbf{$\ell^\infty$ norm:} $\|A\|_{\infty\to\infty} = \max_{1\leq i\leq m}\sum_{j=1}^m |a_{ij}|$ (maximum absolute row sum). \textit{Proof:} For $x$ with $\|x\|_\infty=1$, $|(Ax)_i| = |\sum_j a_{ij}x_j| \leq \sum_j|a_{ij}|$, so $\|Ax\|_\infty \leq \max_i\sum_j|a_{ij}|$, with equality at $x_j = \mathrm{sign}(a_{i^*j})$ for the maximizing row $i^*$.
\end{itemize}

\paragraph{Spectral norm.}
Since $H^\top H=mI_m$, the only eigenvalue of $H^\top H$ is $m$ (with multiplicity $m$). By the spectral norm formula,
\[
\|H\|_{2\to 2} = \sqrt{\lambda_{\max}(H^\top H)} = \sqrt{m}.
\]

\paragraph{$\ell^1$ operator norm.}
Since $|H_{ij}|=1$ for every $i,j$, each column has absolute column sum $\sum_{i=1}^m |H_{ij}| = m$. By the $\ell^1$ norm formula, $\|H\|_{1\to 1}=m$.

\paragraph{$\ell^\infty$ operator norm.}
Similarly, each row has absolute row sum $\sum_{j=1}^m |H_{ij}| = m$. By the $\ell^\infty$ norm formula, $\|H\|_{\infty\to \infty}=m$.

\subsection*{(d) $Hh_i=me_i$}

Let $h_i$ be the $i$th row of $H$, viewed as a column vector. Since $H$ is symmetric, $h_i = H^\top e_i = He_i$. Using part~(b),
\[
Hh_i = H(He_i) = H^2e_i = mI_me_i = me_i. \qedhere
\]
